\section{How much of a menu change can a platform safely make?}\label{sec:protect-main}
Consider a district that has the submitted rankings and two proposed five-course menus. A family already placed under the first policy has a reasonable question about switching to the second: can I keep my place, or get one I prefer? The answer depends on when the district acts.

\paragraph{Before the priority draw}
For each possible draw $s$, clear the old menus to obtain $M^s$. To keep this allocation stable, each pupil must still see their assigned course. A newly shown course must also pass a test: if the pupil prefers it to their assignment, that course must be full with higher-priority holders. For an explicit finite set of draws, protect the union of required courses for each pupil and exclude the union of pairs that fail the test. Fill the remaining slots by score, separately for each pupil. The old menus ensure feasibility. This computes the best additive score under the requirement. In each protected draw, applicant-proposing DA on the new menus gives every pupil an assignment at least as good as $M^s$.

Allowing a fixed fraction of the listed draws to fail changes the optimization. An edit's cost now depends on which draws the other edits already cause to fail. The decision problem is NP-complete even with unit seats, two slots per editable pupil, and one optional substitution per row. The proof encodes graph edges as edits and their endpoints as the two draws in which a pupil loses a seat. A budget of $b$ failed draws can support $\binom b2$ edits exactly when the graph contains a $b$-clique. This hardness concerns an explicitly supplied distribution of orders. It does not establish hardness for the particular independent lottery simulator used below. Appendix~\ref{sec:scenario-boundary} gives the proof and the distinction between a fixed and a draw-dependent baseline.

In the 2022--24 Chilean cohorts we retained the courses shared by the two predictors, then adopted as many remaining rank-informed choices as possible. Protecting all 16 design draws retained 39,293 of 67,909 proposed slot changes. Under the recorded priorities this screen reduced gross exits from 2,088 under the complete switch to 146. Some pupil nevertheless lost a seat in every one of the 192 fresh cohort--lottery draws. These numbers count pupils who lost a placement even when someone else gained one. Increasing the design sample to 1,024 draws kept 29,018 substitutions. At the recorded priorities no original recipient lost a placement, while eight random controls with the same number of edits displaced 894--968 pupils. Yet the screened menus contained 27 \emph{fewer} pupils' eventual schools than the old menus. The full switch had gained 3,590 on the same three cohorts. Could a better choice of edits recover that gain? Even using the April labels in hindsight, the best menu satisfying the same constraints could gain only 278. In 768 further cohort--lottery draws, run after freezing the screened menus, someone lost a seat in 83. Appendix~\ref{sec:lottery-tests} gives the label ceilings and holdout results.

\paragraph{After tentative offers}
Suppose Ana and Ben see only $x$, which gives Ana priority. Ben goes unplaced. Ana prefers the vacant $y$ to $x$. Adding $y$ for Ana can place Ben without changing Ben's menu. At a school with several possible claimants, a move also has to respect people who already saw that school. For each program record the highest priority among applicants who saw it and preferred it to their old offer. A new holder must meet this threshold. Starting at the old matching, find a maximum flow of such moves from unplaced applicants to vacancies. Every moving incumbent gets a course they prefer, and every chain adds one recipient. Implementing any collection of complete chains requires at most one menu substitution per participant. DA on those menus preserves the witness's recipients and weakly improves their assignments. Dropping the program-priority test gives a second flow, an upper bound on what any reassignment can achieve while protecting applicants' welfare (Appendix~\ref{sec:vacancy-chains}).

\begin{table}[ht]\centering\small
\caption{How many additional people can be placed after the rank-informed or Haiku allocation? Direct referrals retain all incumbent contracts. Chains permit preferred moves. The ceiling relaxes program priorities. Deficit is the initial shortfall relative to the unchanged rank-blind or uniform policy. NYC counts pool repeated simulated renters.}\label{tab:design-limits}
\begin{tabular}{@{}lrrrr@{}}\toprule
Starting point & Deficit & Direct referrals & Chains & Ceiling\\\midrule
Chile, four cohorts & 1,787 & 409 & 475 & 649\\
NYC ordinary, 240 markets & 470 & 683 & 983 & 1,353\\
NYC enriched, 240 markets & 898 & 1,926 & 2,493 & 2,844\\\bottomrule
\end{tabular}
\end{table}

In Chile even the ceiling leaves $1,787-649=1,138$ fewer pupils placed than under the unchanged rank-blind menus. Requiring newly shown courses to be in pupils' original top five reduces the ceiling to four. In NYC, the chains recover the unchanged uniform headcount in 226 of 240 ordinary markets and every enriched market. The identity of recipients can differ across starting policies. Applying the same chain algorithm to both rental policies leaves uniform ahead by 128 and 51 placements respectively. The contrast concerns available acceptable moves in these markets, conditional on submitted school ranks and modeled rental preferences. It gives no estimate of leases signed or enrollment caused by a menu.
